\section*{Заключение}
\addcontentsline{toc}{section}{Заключение}

В ходе выполнения лабораторной работы разработано консольное приложение для
преобразования вещественных чисел в формат bfloat16 и обратно. Реализованы ручное
построение двоичной записи целой и дробной частей, нормализация числа, вычисление
смещённой экспоненты, округление семибитной мантиссы и обработка специальных
значений.

Применение структуры с отдельными полями знака, экспоненты и мантиссы позволило
наглядно отразить логическое устройство bfloat16. Результаты проверки показали,
что обратное преобразование восстанавливает ближайшее представимое значение, а
различие с исходным числом определяется ограниченной точностью мантиссы.

\newpage
